
\section{Introduction}

Electroencephalography is routinely affected by ocular activity, and blink contamination remains one of the oldest and most persistent obstacles in EEG analysis. At the same time, spontaneous and task-related blinks are themselves scientifically meaningful. Prior work links blink-derived measures to fatigue, drowsiness, attention, cognitive load, and operational state, and more recent blink-related EEG studies show that blink-linked neural activity can itself index cognitive demands in realistic tasks \citep{croft2000removal,kleifges2017blinker,alyan2023blink}. These dual roles make blink detection a foundational problem rather than a preprocessing detail: researchers need detectors that are accurate enough for artifact handling, but also reliable enough to preserve blink events as analyzable signals in their own right.

Thresholding remains central to this problem because even sophisticated systems frequently begin with amplitude, prominence, duration, or candidate-window gates. Threshold-based designs are attractive in practice because they are interpretable, fast, and deployable in few-channel and real-time environments. That is especially important for frontal wearable EEG and neuroergonomic settings, where channel count, latency, and computational budget are constrained. However, thresholding is also where many detectors fail: a threshold learned from background-dominated samples may drift toward over-detection, while a threshold tuned to one subject or one experiment can collapse under inter-subject variability or under distributions that include eye movements, muscle activity, or epileptiform events \citep{chang2016detection,kleifges2017blinker,agarwal2019blink,tran2021detection,wang2022multidimensional,zhang2023method}.

The central hypothesis of this paper is that threshold selection should be treated as a first-class experimental variable. Instead of asking only which detector wins, we ask which threshold family remains stable when preprocessing, channels, event matching, and datasets are held constant. This paper therefore frames blink detection as a threshold-selection benchmark over interval events. The benchmark explicitly separates signal conditioning, candidate peak extraction, threshold family, region formation, and event matching. It further adds a tightly bounded LLM-assisted ideation module in which the model suggests new threshold derivatives from the literature, but only human-vetted formulas are benchmarked \citep{cao2021unsupervised,saito2015precision,salles2024softed,gottweis2025toward}.

The contributions are fourfold. First, we formalize a threshold taxonomy for blink-region detection in biosignals. Second, we propose a cross-dataset protocol that distinguishes supervised blink benchmarking from out-of-domain robustness testing. Third, we justify overlap-based event scoring with precision, recall, and F1 as the main evaluation regime for sparse blink intervals. Fourth, we introduce an LLM-guided but human-constrained mechanism for proposing interpretable threshold derivatives that can be evaluated under the same benchmark as literature-derived baselines.
